\scanpage{99}
same thing, for the connected components do not change on adding or removing flattened branches), one finds the most general relation between $G$ and a $G'_0$, of the type made explicit in the definition. of ``normalization morphisms of m.graphs''.

One defines in an obvious way the composites of normalization morphisms. The composite of strict \unsure{normalization} morphisms is strict. The m.graphs with strict normalization morphisms form a category, whose isomorphisms are the ordinary isomorphisms of m.graphs. The morphisms of this category are monomorphisms. When one restricts to strict \inserted{normalization} morphisms, one sees that (as for generization morphisms) the set of strict sub-m.graphs of an m.graph \inserted{$G$} is an ordered set \deleted{\unsure{iso}}\inserted{anti}morphic to the set of subsets of $A(G)$, by associating to every subset $C\subset A(\unsure{G|C})$
\margin{98}
